% =====================================================================
% CAP. 11 - NIVEL 2 - FONTES DEPENDENTES
% =====================================================================
\blocoaplicacao

\begin{questao}[2]{}
O circuito abaixo parte com o capacitor descarregado. Em $t=0$, a fonte de corrente
independente aplica um degrau de $\SI{3}{\milli\ampere}$. A fonte dependente é uma
VCCS de valor $g v_C$, com $g=\SI{0.25}{\milli\ensuremath{\mathrm{S}}}$. Determine a constante de
tempo, o valor final de $v_C$ e $v_C(\SI{0.10}{\second})$.

\circuitoqq{%
\begin{circuitikz}[scale=0.96,transform shape,line width=0.8pt,every node/.append style={font=\scriptsize}]
\draw (0,0) to[american current source] (0,3) -- (7.0,3);
\node[anchor=east] at (-0.25,1.5) {$3u(t)\,\mathrm{mA}$};
\draw (2.0,3) to[R] (2.0,0);
\node[anchor=east] at (1.75,1.5) {$R=\SI{4}{\kilo\ohm}$};
\draw (4.5,3) to[C,v^>={$v_C$}] (4.5,0);
\node[anchor=east] at (4.25,1.5) {$C=\SI{50}{\micro\farad}$};
\draw (7.0,3) to[american controlled current source] (7.0,0) -- (0,0);
\node[anchor=west] at (7.25,1.5) {$g\,v_C$};
\end{circuitikz}}{Circuito RC com fonte de corrente controlada por tensão (VCCS).}

\resposta{$R_{\mathrm{eq}}=\SI{2}{\kilo\ohm}$; $\tau=\SI{0.10}{\second}$;
$v_C(\infty)=\SI{6}{\volt}$; $v_C(0{,}10\,\mathrm{s})\approx\SI{3.79}{\volt}$}
\resolucao{A fonte dependente não é suprimida. A condutância total vista pelo
capacitor é
\[
G_{\mathrm{eq}}=\frac1R+g=\SI{0.25}{\milli\ensuremath{\mathrm{S}}}+
\SI{0.25}{\milli\ensuremath{\mathrm{S}}}=\SI{0.50}{\milli\ensuremath{\mathrm{S}}},
\]
logo $R_{\mathrm{eq}}=1/G_{\mathrm{eq}}=\SI{2}{\kilo\ohm}$ e
$\tau=R_{\mathrm{eq}}C=\SI{0.10}{\second}$. Em regime permanente,
$v_C=I_s/G_{\mathrm{eq}}=\SI{6}{\volt}$. Portanto
$v_C(t)=6(1-e^{-t/0{,}1})$ e, em $t=\SI{0.10}{\second}$,
$v_C\approx\SI{3.79}{\volt}$.}
\end{questao}

\begin{questao}[2]{}
O capacitor do circuito natural abaixo está inicialmente a $\SI{10}{\volt}$.
A fonte dependente fornece a corrente $g v_C$, com
$g=\SI{0.5}{\milli\ensuremath{\mathrm{S}}}$. Determine $R_{\mathrm{eq}}$, $\tau$ e $v_C(t)$.

\circuitoqq{%
\begin{circuitikz}[scale=0.98,transform shape,line width=0.8pt,every node/.append style={font=\scriptsize}]
\draw (0,0) -- (7.0,0);
\draw (0,3) -- (7.0,3);
\draw (1.8,3) to[R] (1.8,0);
\node[anchor=east] at (1.55,1.5) {$R=\SI{2}{\kilo\ohm}$};
\draw (4.2,3) to[C,v^>={$v_C$}] (4.2,0);
\node[anchor=east] at (3.95,1.5) {$C=\SI{100}{\micro\farad}$};
\draw (7.0,3) to[american controlled current source] (7.0,0);
\node[anchor=west] at (7.25,1.5) {$g\,v_C$};
\end{circuitikz}}{Resposta natural RC com uma VCCS em paralelo.}

\resposta{$R_{\mathrm{eq}}=\SI{1}{\kilo\ohm}$; $\tau=\SI{0.10}{\second}$;
$v_C(t)=10e^{-10t}\ \si{\volt}$}
\resolucao{Como a fonte dependente permanece ativa,
$G_{\mathrm{eq}}=1/R+g=\SI{0.5}{\milli\ensuremath{\mathrm{S}}}+
\SI{0.5}{\milli\ensuremath{\mathrm{S}}}=\SI{1}{\milli\ensuremath{\mathrm{S}}}$. Logo
$R_{\mathrm{eq}}=\SI{1}{\kilo\ohm}$ e
$\tau=R_{\mathrm{eq}}C=\SI{0.10}{\second}$. Como não há fonte independente e
o valor final é zero, $v_C(t)=10e^{-t/\tau}=10e^{-10t}\,\si{\volt}$.}
\end{questao}

\begin{questao}[2]{}
No circuito abaixo, o indutor possui $i_L(0^+)=\SI{2}{\ampere}$. A fonte dependente
é uma CCCS que injeta para cima a corrente $\beta i_x$, em que
$i_x=v/R$ e $\beta=0{,}5$. Determine a resistência equivalente vista pelo indutor,
a constante de tempo e $i_L(t)$.

\circuitoqq{%
\begin{circuitikz}[scale=0.98,transform shape,line width=0.8pt,every node/.append style={font=\scriptsize}]
\draw (0,0) -- (7.2,0);
\draw (0,3) -- (7.2,3);
\draw (1.8,3) to[R,i>^={$i_x$}] (1.8,0);
\node[anchor=east] at (1.55,1.5) {$R=\SI{20}{\ohm}$};
\draw (4.2,3) to[L,i>^={$i_L$}] (4.2,0);
\node[anchor=east] at (3.95,1.5) {$L=\SI{200}{\milli\henry}$};
\draw (7.2,0) to[american controlled current source] (7.2,3);
\node[anchor=west] at (7.45,1.5) {$\beta\,i_x$};
\end{circuitikz}}{Resposta natural RL com fonte de corrente controlada por corrente (CCCS).}

\resposta{$R_{\mathrm{eq}}=\SI{40}{\ohm}$; $\tau=\SI{5}{\milli\second}$;
$i_L(t)=2e^{-200t}\ \si{\ampere}$}
\resolucao{Aplicando uma tensão-teste $v_x$ aos terminais do indutor, a corrente
no resistor é $i_x=v_x/R$. A fonte dependente injeta $\beta i_x$ no nó superior,
logo a corrente fornecida pela fonte-teste é
\[
i_t=i_x-\beta i_x=(1-\beta)\frac{v_x}{R}.
\]
Portanto
$R_{\mathrm{eq}}=v_x/i_t=R/(1-\beta)=20/0{,}5=\SI{40}{\ohm}$.
Assim $\tau=L/R_{\mathrm{eq}}=0{,}2/40=\SI{5}{\milli\second}$ e
$i_L(t)=2e^{-t/\tau}=2e^{-200t}\,\si{\ampere}$.}
\end{questao}

\gabarito[Nível 2 --- fontes dependentes]
